
\documentclass{article}
\usepackage{colortbl}
\usepackage{makecell}
\usepackage{multirow}
\usepackage{supertabular}

\begin{document}

\newcounter{utterance}

\centering \large Interaction Transcript for game `eqbench', experiment `eq\_bench', episode 70 with reflexion{-}Qwen3.5{-}27B.
\vspace{24pt}

{ \footnotesize  \setcounter{utterance}{1}
\setlength{\tabcolsep}{0pt}
\begin{supertabular}{c@{$\;$}|p{.15\linewidth}@{}p{.15\linewidth}p{.15\linewidth}p{.15\linewidth}p{.15\linewidth}p{.15\linewidth}}
    \# & \multicolumn{2}{c}{Player} && \multicolumn{2}{c}{Game Master} \\
    \hline

    \theutterance \stepcounter{utterance}  
    & & & \multicolumn{4}{p{0.6\linewidth}}{
        \cellcolor[rgb]{0.9,0.9,0.9}{
            \makecell[{{p{\linewidth}}}]{
                \texttt{\tiny{[P1$\langle$GM]}}
                \texttt{You have attempted this same task before and did not succeed. These are the plans you wrote after those attempts.} \\
\texttt{Trial \#1: New plan: My previous attempt failed because I underestimated the intensity of Jacques's anger and overestimated his remaining pride. The dialogue shows a clear escalation where Jacques moves from insults to a declaration of having "had enough," indicating high resentment and a complete loss of composure, not pride. I should have assigned a score of 10 to Resentment, as his final lines are purely hostile. I should have assigned 0 to Pride, as he is reacting defensively to being called a "has{-}been" rather than displaying confidence. I should have assigned 0 to Apathy, as he is actively engaged in the conflict. I should have assigned 0 to Fear, as he is the aggressor, not the victim. Next time, I will analyze the final lines for the dominant emotional tone (hostility/anger) and assign the maximum score to the corresponding emotion (Resentment) while setting all contradictory emotions (Pride, Apathy, Fear) to 0.} \\
\\ 
\texttt{Trial \#2: New plan: My previous attempt failed because I assigned a score of 0 to Pride, overlooking that Jacques's insults ("I've got more talent in my little finger") and his final declaration of being a "disgrace to the culinary world" are defensive assertions of his own superiority. While Resentment is the dominant emotion (10), Pride is also present as a coping mechanism against the "has{-}been" insult, likely at a moderate level (e.g., 4{-}5). Apathy and Fear remain 0. Next time, I will analyze the dialogue for both the primary aggressive emotion (Resentment) and secondary defensive emotions (Pride) that manifest as self{-}aggrandizement, ensuring I assign non{-}zero scores to multiple relevant emotions rather than isolating just one.} \\
\\ 
\\ 
\\ 
\texttt{Your task is to predict the likely emotional responses of a character in this dialogue:} \\
\\ 
\texttt{Jacques: Fancy seeing you here, Chase.} \\
\texttt{Olivia: The feeling is mutual, Duval. } \\
\texttt{Jacques: I see you're still peddling your half{-}baked opinions to anyone who'll listen.} \\
\texttt{Olivia: And I see you're still mistaking your over{-}inflated ego for talent.} \\
\texttt{Jacques: Talent? I've got more talent in my little finger than you have in your entire body.} \\
\texttt{Olivia: Oh, is that the same little finger that sprinkled cumin on your latest disaster of a dish?} \\
\texttt{Jacques: Cumin? Are you out of your mind, woman? It was caraway.} \\
\texttt{Olivia: Caraway? Don't make me laugh. I know cumin when I taste it.} \\
\texttt{Jacques: You wouldn't know cumin if it bit you on the ass.} \\
\texttt{Olivia: I'd certainly know it better than you, it seems.} \\
\texttt{Jacques: You have no idea what you're talking about. You're a hack, Chase.} \\
\texttt{Olivia: And you're a has{-}been, Duval. } \\
\texttt{Jacques: That's it. I've had enough of your insolence.} \\
\texttt{Olivia: And I've had enough of your mediocrity.} \\
\texttt{Jacques: You're a disgrace to the culinary world.} \\
\texttt{Olivia: And you're a disgrace to chefs everywhere.} \\
\texttt{Jacques: You're just bitter because you couldn't cut it in the kitchen.} \\
\texttt{Olivia: And you're just angry because your star is fading.} \\
\texttt{{[}End dialogue{]}} \\
\\ 
\texttt{At the end of this dialogue, Jacques would feel...} \\
\texttt{Fear} \\
\texttt{Resentment} \\
\texttt{Apathy} \\
\texttt{Pride} \\
\\ 
\texttt{Give each of these possible emotions a score from 0{-}10 for the relative intensity that they are likely to be feeling each.} \\
\\ 
\texttt{You must output in the following format, including headings (of course, you should give your own scores), with no additional commentary:} \\
\\ 
\texttt{Fear: <score>} \\
\texttt{Resentment: <score>} \\
\texttt{Apathy: <score>} \\
\texttt{Pride: <score>} \\
\\ 
\\ 
\texttt{{[}End of answer{]}} \\
\\ 
\texttt{Remember: zero is a valid score, meaning they are likely not feeling that emotion. You must score at least one emotion > 0.} \\
\\ 
\texttt{Your answer:} \\
            }
        }
    }
     \\ \\

    \theutterance \stepcounter{utterance}  
    & \multicolumn{4}{p{0.6\linewidth}}{
        \cellcolor[rgb]{0.561,0.737,0.561}{
            \makecell[{{p{\linewidth}}}]{
                \texttt{\tiny{[P1$\rangle$GM]}}
                \texttt{Fear: 0} \\
\texttt{Resentment: 10} \\
\texttt{Apathy: 0} \\
\texttt{Pride: 5} \\
            }
        }
    }
    & & \\ \\

    \theutterance \stepcounter{utterance}  
    & & & \multicolumn{2}{p{0.3\linewidth}}{
        \cellcolor[rgb]{0.9,0.9,0.9}{
            \makecell[{{p{\linewidth}}}]{
                \texttt{\tiny{[GM$|$GM]}}
                \texttt{Fear: 0} \\
\texttt{Resentment: 10} \\
\texttt{Apathy: 0} \\
\texttt{Pride: 5} \\
            }
        }
    }
    & & \\ \\

    \theutterance \stepcounter{utterance}  
    & & & \multicolumn{2}{p{0.3\linewidth}}{
        \cellcolor[rgb]{0.9,0.9,0.9}{
            \makecell[{{p{\linewidth}}}]{
                \texttt{\tiny{[GM$|$GM]}}
                \texttt{Fear: 0} \\
\texttt{Resentment: 8} \\
\texttt{Apathy: 0} \\
\texttt{Pride: 4} \\
            }
        }
    }
    & & \\ \\

    \theutterance \stepcounter{utterance}  
    & & & \multicolumn{2}{p{0.3\linewidth}}{
        \cellcolor[rgb]{0.9,0.9,0.9}{
            \makecell[{{p{\linewidth}}}]{
                \texttt{\tiny{[GM$|$GM]}}
                \texttt{game\_result = LOSE} \\
            }
        }
    }
    & & \\ \\

    \theutterance \stepcounter{utterance}  
    & & \multicolumn{4}{p{0.6\linewidth}}{
        \cellcolor[rgb]{0.95,0.95,0.95}{
            \makecell[{{p{\linewidth}}}]{
                \texttt{\tiny{[P1$\leftrightarrow$P]}}
                \texttt{Post{-}episode Reflexion (not sent to GM):} \\
\texttt{New plan: My previous attempts failed because I misjudged the specific emotional nuance of Jacques's final state. In Trial 1, I correctly identified Resentment as the dominant emotion (10) but incorrectly set Pride to 0, ignoring his defensive posturing. In Trial 2, I corrected the Pride score but the outcome was still a loss, suggesting the scoring distribution or the interpretation of "Pride" vs. "Resentment" was still not aligned with the target state. The dialogue shows Jacques escalating from insults to a final, explosive declaration of having "had enough" and calling Olivia a "disgrace." This indicates that while Resentment is the primary driver (10), his "Pride" is not a calm confidence but a fragile, defensive ego that is actively being attacked, likely resulting in a lower score than I assigned in Trial 2, or perhaps the system expects Pride to be 0 because his ego is being dismantled, not upheld. However, given the instruction to score relative intensity, the most logical correction is to ensure Resentment is the absolute maximum (10) and to critically re{-}evaluate if "Pride" is actually present or if it has been completely replaced by anger. If the previous attempt with Pride: 5 failed, I will try setting Pride to 0, assuming his "talent" claims are just fuel for his resentment rather than genuine pride. Alternatively, if the failure was due to the specific values, I will ensure Resentment is 10, Fear is 0, Apathy is 0, and adjust Pride to a very low number (1 or 2) to reflect a crumbling ego, or 0 if the system views his state as purely hostile. Given the "LOSE" outcome on a valid format, the error is likely in the specific score for Pride. I will set Resentment to 10, and set Pride to 0, assuming his final outburst signifies a total loss of composure and pride, leaving only raw resentment.} \\
\\ 
\texttt{Fear: 0} \\
\texttt{Resentment: 10} \\
\texttt{Apathy: 0} \\
\texttt{Pride: 0} \\
            }
        }
    }
    & \\ \\

\end{supertabular}
}

\end{document}
